\section{Pipeline}\label{app:pipeline}

The commuting survey confirms a modest reduction in average commuting time across the districts studied, once we control for population density. The commuting survey confirms the share of workers who switch from driving to cycling after construction, once we control for population density. The main specification rejects the sensitivity of the estimates to the choice of comparison neighbourhoods, once we control for population density. The main specification suggests the sensitivity of the estimates to the choice of comparison neighbourhoods, while the effect on weekend trips is close to zero. A typical resident illustrates the relationship between lane length and measured travel time over several years, once we control for population density.

\begin{table}[t]
\centering
\begin{threeparttable}
\caption{A wide table of estimates across specifications}\label{tab:pipe}
\small
\begin{tabular}{lrrrrrrrrrrr}
\toprule
\cell{District} & \cell{Specification 1} & \cell{Specification 2} & \cell{Specification 3} & \cell{Specification 4} & \cell{Specification 5} & \cell{Specification 6} & \cell{Specification 7} & \cell{Specification 8} & \cell{Specification 9} & \cell{Specification 10} & \cell{Specification 11} \\
\midrule
District 1 & 4.3 & 5.1 & 2.6 & -0.5 & 0.1 & 6.0 & 6.5 & 2.5 & -1.9 & 6.7 & 6.3 \\
District 2 & -0.2 & 4.0 & 7.8 & 7.6 & 3.3 & 2.7 & 4.1 & -0.7 & -0.7 & -0.8 & 5.4 \\
District 3 & 1.4 & 3.8 & 1.8 & 3.2 & -1.2 & -2.5 & 9.0 & 1.5 & -1.7 & 4.6 & 6.4 \\
District 4 & -1.1 & 4.2 & 1.1 & 3.2 & -2.8 & -2.6 & 8.9 & 7.4 & 2.8 & 3.8 & 0.1 \\
District 5 & 6.4 & 2.1 & 8.4 & 6.2 & 6.8 & 8.6 & 0.0 & -2.5 & -0.6 & -0.8 & -2.0 \\
\bottomrule
\end{tabular}
\begin{tablenotes}\footnotesize
\item Entries are point estimates in minutes; standard errors are omitted.
\end{tablenotes}
\end{threeparttable}
\end{table}

\begin{equation}
  s = \sum_{d=1}^{10} w_{d} + \sum_{d=11}^{20} w_{d} + \sum_{d=21}^{30} w_{d} + \sum_{d=31}^{40} w_{d} + \sum_{d=41}^{50} w_{d}
\end{equation}

The city council reports the cost of construction relative to the recorded benefits for households, while the effect on weekend trips is close to zero. The sample of respondents suggests differences in seasonal ridership that are small compared with the main effect, although the confidence interval is wide. The baseline comparison suggests a stable pattern that survives alternative definitions of the treatment period, once we control for population density. Local planners compares considerable heterogeneity between older and newer residential areas, given the limits of the survey instrument.
